% GENERATED FILE. Edit canonical lessons, then run npm run generate:books.
% canonical-source-hash: fnv1a64:387e20286d2b7b9e
% canonical-lessons: PA-C82-bharia, PA-C82-khali, PA-C82-saf, PA-C82-mittha, PA-C82-mota

\chapter{Full, Empty, Clean, Sweet, Fat}
\label{ch:pa-full-empty-clean-sweet-fat}
\hlchaptermodality{82}

\begin{chapteropening}
\textbf{By the end of this chapter:} \emph{I can say full, empty, clean, sweet and fat, thick.}

Complete the last lesson of chapter 82: i can say full, empty, clean, sweet and fat, thick.
\end{chapteropening}

\section[bhariā]{\pa{ਭਰਿਆ} (bhariā) — full}

\label{lesson:PA-C82-bharia}



\begin{quote}
Before the new one: say the Punjabi for high, then the Punjabi for low.
\end{quote}

\subsection*{You'll want to know: \pa{ਭਰਿਆ}}
\textbf{\pa{ਭਰਿਆ}} — \emph{bhariā} — \textquotedblleft{}full\textquotedblright{}.

\begin{grammarlens}[title={What you've built}]
One more word for this chapter.
\end{grammarlens}

\subsection*{Your turn}
\begin{itemize}
\raggedright
  \item \emph{Say it:} \emph{bhariā}
  \item \emph{Say it:} \emph{bhariā}, once more
  \item \emph{Say it:} say \emph{bhariā}
  \item \emph{Recall:} say the Punjabi for high, then the Punjabi for low, then say \emph{bhariā} again
\end{itemize}

\begin{tcolorbox}[breakable,colback=teal!4,colframe=teal!35!black,title={Before you move on}]
What does \pa{ਭਰਿਆ} mean? (\textquotedblleft{}Full\textquotedblright{}.) Say it once more.
\end{tcolorbox}
\section[khālī]{\pa{ਖ਼ਾਲੀ} (khālī) — empty}

\label{lesson:PA-C82-khali}



\begin{quote}
Before the new one: say the Punjabi for low, then the Punjabi for full.
\end{quote}

\subsection*{You'll want to know: \pa{ਖ਼ਾਲੀ}}
\textbf{\pa{ਖ਼ਾਲੀ}} — \emph{khālī} — \textquotedblleft{}empty\textquotedblright{}.

\begin{grammarlens}[title={What you've built}]
One more word for this chapter.
\end{grammarlens}

\subsection*{Your turn}
\begin{itemize}
\raggedright
  \item \emph{Say it:} \emph{khālī}
  \item \emph{Say it:} \emph{khālī}, once more
  \item \emph{Say it:} say \emph{khālī}
  \item \emph{Recall:} say the Punjabi for low, then the Punjabi for full, then say \emph{khālī} again
\end{itemize}

\begin{tcolorbox}[breakable,colback=teal!4,colframe=teal!35!black,title={Before you move on}]
What does \pa{ਖ਼ਾਲੀ} mean? (\textquotedblleft{}Empty\textquotedblright{}.) Say it once more.
\end{tcolorbox}
\section[sāf]{\pa{ਸਾਫ਼} (sāf) — clean}

\label{lesson:PA-C82-saf}



\begin{quote}
Before the new one: say the Punjabi for full, then the Punjabi for empty.
\end{quote}

\subsection*{You'll want to know: \pa{ਸਾਫ਼}}
\textbf{\pa{ਸਾਫ਼}} — \emph{sāf} — \textquotedblleft{}clean\textquotedblright{}.

\begin{grammarlens}[title={What you've built}]
One more word for this chapter.
\end{grammarlens}

\subsection*{Your turn}
\begin{itemize}
\raggedright
  \item \emph{Say it:} \emph{sāf}
  \item \emph{Say it:} \emph{sāf}, once more
  \item \emph{Say it:} say \emph{sāf}
  \item \emph{Recall:} say the Punjabi for full, then the Punjabi for empty, then say \emph{sāf} again
\end{itemize}

\begin{tcolorbox}[breakable,colback=teal!4,colframe=teal!35!black,title={Before you move on}]
What does \pa{ਸਾਫ਼} mean? (\textquotedblleft{}Clean\textquotedblright{}.) Say it once more.
\end{tcolorbox}
\section[miṭṭhā]{\pa{ਮਿੱਠਾ} (miṭṭhā) — sweet}

\label{lesson:PA-C82-mittha}



\begin{quote}
Before the new one: say the Punjabi for empty, then the Punjabi for clean.
\end{quote}

\subsection*{You'll want to know: \pa{ਮਿੱਠਾ}}
\textbf{\pa{ਮਿੱਠਾ}} — \emph{miṭṭhā} — \textquotedblleft{}sweet\textquotedblright{}.

\begin{grammarlens}[title={What you've built}]
One more word for this chapter.
\end{grammarlens}

\subsection*{Your turn}
\begin{itemize}
\raggedright
  \item \emph{Say it:} \emph{miṭṭhā}
  \item \emph{Say it:} \emph{miṭṭhā}, once more
  \item \emph{Say it:} say \emph{miṭṭhā}
  \item \emph{Recall:} say the Punjabi for empty, then the Punjabi for clean, then say \emph{miṭṭhā} again
\end{itemize}

\begin{tcolorbox}[breakable,colback=teal!4,colframe=teal!35!black,title={Before you move on}]
What does \pa{ਮਿੱਠਾ} mean? (\textquotedblleft{}Sweet\textquotedblright{}.) Say it once more.
\end{tcolorbox}
\section[moṭā]{\pa{ਮੋਟਾ} (moṭā) — fat, thick}

\label{lesson:PA-C82-mota}



\begin{quote}
Before the new one: say the Punjabi for clean, then the Punjabi for sweet.
\end{quote}

\subsection*{You'll want to know: \pa{ਮੋਟਾ}}
\textbf{\pa{ਮੋਟਾ}} — \emph{moṭā} — \textquotedblleft{}fat, thick\textquotedblright{}.

\begin{grammarlens}[title={What you've built}]
One more word for this chapter.
\end{grammarlens}

\subsection*{Your turn}
\begin{itemize}
\raggedright
  \item \emph{Say it:} \emph{moṭā}
  \item \emph{Say it:} \emph{moṭā}, once more
  \item \emph{Say it:} say \emph{moṭā}
  \item \emph{Recall:} say the Punjabi for clean, then the Punjabi for sweet, then say \emph{moṭā} again
\end{itemize}

\begin{tcolorbox}[breakable,colback=teal!4,colframe=teal!35!black,title={Before you move on}]
What does \pa{ਮੋਟਾ} mean? (\textquotedblleft{}Fat, thick\textquotedblright{}.) Say it once more.
\end{tcolorbox}
